% ============================================================
%  CHAPTER 3 — Methodology
%  Typical length: 20–30 pages
%  Cover: theoretical framework, model/dataset description,
%         simulation or experimental design, performance metrics.
% ============================================================

% Section are made as examples to show how to write — you have to make sections/subsections as per your need.


\chapter{Materials and Methods}
\label{chap:methods}

[Section are made as examples — make your own sections/subsections headings as per your need.]

% TODO: Brief overview paragraph — what does this chapter cover?

% ---- 3.1 Theoretical Framework -----------------------------
\section{Theoretical Framework}
\label{sec:theory}

% TODO: Describe the mathematical or physical model underpinning
% your work. Derive or state the key equations.
% Reference appendices for lengthy derivations.

[Write the theoretical framework here.]

% ---- 3.2 Model / Data Description --------------------------
\section{Model Description}
\label{sec:model}

% TODO: Describe your simulation model, dataset, or experimental
% setup in enough detail that the work could be reproduced.

[Describe your model or dataset here.]

% Example of a table — 3 columns, 4 data rows.
% \toprule / \midrule / \bottomrule require the booktabs package (already loaded).
% Replace column headers and data with your own values.
% Use \label{tab:your_label} and cross-reference with Table~\ref{tab:your_label}.

\begin{table}[htbp]
  \centering
  \caption{Example table: model parameters used in this study.}
  \label{tab:model_parameters}
  \begin{tabular}{lcc}
    \toprule
    \textbf{Parameter} & \textbf{Value} & \textbf{Unit} \\
    \midrule
    Parameter 1        & 100            & m             \\
    Parameter 2        & 3.5            & km/s          \\
    Parameter 3        & 50             & Hz            \\
    Parameter 4        & 0.01           & --            \\
    \bottomrule
  \end{tabular}
\end{table}

% ---- 3.3 Methods / Workflow --------------------------------
\section{Methods}
\label{sec:workflow}

% TODO: Step-by-step description of your methodology.
% A flowchart figure here is often very effective.

[Describe your workflow and processing steps here.]

\subsection{[Step / Method 1]}

[Write here.]

\subsection{[Step / Method 2]}

[Write here.]

% ---- 3.4 Performance Metrics -------------------------------
\section{Performance Metrics}
\label{sec:metrics}

% TODO: Define each quantitative metric you use to evaluate results.
% Give the mathematical definition and explain what it measures.

[Define your evaluation metrics here.]


% Uncomment the below lines (if you have used AI for the contents of this chapter) to include an AI use disclosure at the end of the chapter. You can modify the wording as needed, but be transparent about how you used AI tools in writing this chapter.

% \vspace{1.0cm}
% \noindent\rule{\textwidth}{0.4pt}
% \vspace{0.2cm}

% \noindent\textit{Note on AI assistance:} Throughout this chapter, an AI language model was used as a collaborative tool for [describe how — e.g., brainstorming model design choices, methodology decisions, metric selection]. All methodological decisions, implementations, and interpretations are the author's own.
